%-*-latex-*-
%-*-latex-*-
\newcommand\COURSE{ciss245}
\newcommand\ASSESSMENT{q8001}
\newcommand\ASSESSMENTTYPE{Quiz}
\newcommand\POINTS{\textwhite{xxx/xxx}}

\input{myquizpreamble}
\input{yliow}
\input{\COURSE}
%-*-latex-*-

%-*-latex-*-

%-*-latex-*-
\renewcommand\TITLE{\ASSESSMENTTYPE \ \ASSESSMENT}

\renewcommand\EMAIL{}
\renewcommand\AUTHOR{eagraham1@cougars.ccis.edu}

\textwidth=6in
\begin{document}
\topmatterthree

%------------------------------------------------------------------------------
\nextq
In the code below, there are common member variables in
\verb!TeslaOptimusRobot!
and
\verb!RoombaRobot!.
The goal is to clean up the repetitive code by moving the common
parts of these two classes into a parent class.
\begin{enumerate}[nosep]
\li Complete the \verb!Robot! class with the common member variables of
\verb!TeslaOptimusRobot! and \verb!Roomba! in its private section.
\li In the public section, write constructor for \verb!Robot! to initialize these member variables
and write public constant methods \verb!Robot::get_name()!,
\verb!Robot::get_x()!, and \verb!Robot::get_y()!.
\li
Make \verb!TeslaOptimusRobot! and \verb!RoombaRobot! subclasses of \verb!Robot!
using public inheritance.
Modify \verb!TeslaOptimusRobot! and \verb!RoombaRobot! accordingly.
\li There should be no change to \verb!main()!
and the completed code should run the same as before.
\end{enumerate}
\ANSWER
\begin{answercode}
#include <iostream>
#include <string>

class Robot
{
public:
Robot(std::string name, int x, int y)
: name_(name), x_(x), y_(y)
{}
std::string get_name() const
{
  return name_;
}
int get_x() const
{
  return x_;
}
int get_y() const
{
  return y_;
}
private:
std::string name_;
int x_, y_;
};

class TeslaOptimusRobot: public Robot
{
public:
    TeslaOptimusRobot(std::string name, int x, int y, int num_legs)
        : Robot(name, x, y), num_legs_(num_legs)
    {}
    void print() const
    {
        std::cout << "<TeslaOptimusRobot "
                  << "name:" << name_ << ' '
                  << "x:" << get_x() << ' '
                  << "y:" << get_y() << ' '
                  << "num_legs:" << num_legs_
                  << ">\n";
    }
    void talk() const
    {
        std::cout << "<TeslaOptimusRobot name:"
                  << get_name()
                  << "> says \"hello world\"\n";
    }
private:
    int num_legs_;
};

class RoombaRobot: public Robot
{
public:
    RoombaRobot(std::string name, int x, int y, int num_wheels)
        : Robot(name, x, y), num_wheels_(num_wheels)
    {}
    void print() const
    {
        std::cout << "<RoombaRobot "
                  << "name:" << get_name() << ' '
                  << "x:" << get_x() << ' '
                  << "y:" << get_y() << ' '
                  << "num_wheels:" << num_wheels_
                  << ">\n";
    }
    void talk() const
    {
        std::cout << "<RoombaRobot name:"
                  << get_name()
                  << "> says \"whirrrr\"\n";
    }
private:
    int num_wheels_;
};

int main()
{
    TeslaOptimusRobot arnold("arnold", 0, 0, 2);
    RoombaRobot neo("i_am_the_chosen_one", 5, 0, 3);

    arnold.print();
    neo.print();

    arnold.talk();
    neo.talk();

    return 0;
}
\end{answercode}

%------------------------------------------------------------------------------
%-*-latex-*-
\newpage
\input{instructions}
\end{document}

